\chapter{Conclusion}
\label{chap:Conclusion}
This project successfully designed and documented the software architecture for the Quiz Web App.
The architectural design, based on a three-tier architecture (presentation, application, and data tiers), provides a scalable and maintainable solution.
The presentation tier, developed using Angular, offers a responsive and interactive user interface.
The application tier, built with Node.js and Express.js, handles business logic and API requests.
The data tier utilizes MongoDB for flexible and efficient data storage.

Key architectural decisions, such as the choice of technologies and the separation of concerns, were driven by the functional and non-functional requirements outlined in the SRS.
The system modeling, including use case diagrams, activity diagrams, sequence diagrams, and class diagrams, provided a clear understanding of the system's behavior and structure.
The requirements definition phase ensured that all stakeholder needs were captured and addressed.

The integration of an external AI API (e.g., OpenAI) for features like question generation assistance enhances the application's capabilities and user experience.
Deployment strategies, including local and Docker-based deployment, were also considered, ensuring ease of setup and portability.
The project utilized tools like \LaTeX{} for documentation, Git for version control, and Visual Studio Code as the primary IDE, which facilitated an efficient development and documentation process.

Future work could involve expanding the feature set, such as incorporating more advanced AI-driven functionalities, supporting a wider range of question types, or enhancing the analytics and reporting capabilities for administrators.
Further refinement of the security measures and performance optimization could also be undertaken as the user base grows.

In summary, the architectural framework established in this project provides a solid foundation for the development and future evolution of the Quiz Web App, meeting the initial objectives of creating a robust, user-friendly, and feature-rich quizzing platform.
